\documentclass[11pt]{article}

\usepackage[margin=1in]{geometry}
\usepackage{amsmath,amssymb,amsthm,mathtools}
\usepackage{hyperref}
\usepackage{enumitem}

\hypersetup{colorlinks=true,linkcolor=blue,urlcolor=blue,citecolor=blue}
\emergencystretch=2em

\newtheorem{proposition}{Proposition}
\newcommand{\C}{\mathbb C}

\title{An Isometric Operator-Algebra Action Which Is Not Completely Isometric}
\author{Run fa\_banach\_001, lane 7}
\date{August 11, 2026}

\begin{document}
\maketitle

\section{Question}

Wang and Zhu ask in Remark 2.4(i) of arXiv:2305.03933 whether their
$p$-nuclearity theorem for crossed products remains true when a
$p$-completely isometric action is replaced by an isometric action.  They
also state the subsidiary question whether an isometric action is necessarily
$p$-completely isometric.  The latter has a negative answer as written,
already when $p=2$.

\section{Construction}

Let $e_{ij}$ denote the standard matrix units in $M_3$, acting on
$H=\ell_2^3$.  Put
\[
 C_1=e_{21},\qquad C_2=e_{31},\qquad
 R_1=e_{12}=C_1^{\mathsf t},\qquad
 R_2=e_{13}=C_2^{\mathsf t},
\]
and define
\[
 B_C=\operatorname{span}\{I,C_1,C_2\},\qquad
 B_R=\operatorname{span}\{I,R_1,R_2\}.
\]
All products $C_jC_k$ and $R_jR_k$ vanish.  Hence $B_C$ and $B_R$ are
unital commutative operator algebras.  Transpose gives a complex-linear
algebra isomorphism
\[
 \theta:B_C\longrightarrow B_R,\qquad \theta(T)=T^{\mathsf t}.
\]
For the Hilbert-space operator norm, $\|T^{\mathsf t}\|=\|T\|$; equivalently,
$T$ and $T^{\mathsf t}$ have the same singular values.  Thus $\theta$ is an
isometry at level one.

Represent
\[
 A=B_C\oplus B_R
 \subset B(H\oplus_2H)
\]
as a block-diagonal unital operator algebra, and define
\[
 \alpha(b,c)=\big(\theta^{-1}(c),\theta(b)\big).
\]
The direct-sum norm is the maximum of the two component norms.  Therefore
$\alpha$ is an isometric algebra automorphism and $\alpha^2=\mathrm{id}$.
It generates an isometric action of the finite amenable group $\mathbb Z_2$.

\begin{proposition}
The action generated by $\alpha$ is not $2$-completely isometric.
\end{proposition}

\begin{proof}
In $M_2(A)$ consider
\[
 X=
 \begin{bmatrix}
  (C_1,0)&0\\
  (C_2,0)&0
 \end{bmatrix}.
\]
The nonzero direct-sum component of $X$ is the operator
\[
 \mathsf C=
 \begin{bmatrix}C_1&0\\C_2&0\end{bmatrix}
 \quad\text{on }H\oplus_2H.
\]
If $\xi=(\xi_1,\xi_2)$ and $\xi_1=(z_0,z_1,z_2)$, then
\[
 \mathsf C\xi=(C_1\xi_1,C_2\xi_1),
 \qquad \|\mathsf C\xi\|^2=2|z_0|^2.
\]
Consequently $\|X\|=\|\mathsf C\|=\sqrt2$, with equality attained at
$\xi_1=(1,0,0)$ and $\xi_2=0$.

Entrywise application of $\alpha$ gives
\[
 \alpha^{(2)}(X)=
 \begin{bmatrix}
  (0,R_1)&0\\
  (0,R_2)&0
 \end{bmatrix}.
\]
Its nonzero component is
\[
 \mathsf R=
 \begin{bmatrix}R_1&0\\R_2&0\end{bmatrix}.
\]
For the same vector $\xi$,
\[
 \mathsf R\xi=(R_1\xi_1,R_2\xi_1),
 \qquad \|\mathsf R\xi\|^2=|z_1|^2+|z_2|^2.
\]
Thus $\|\alpha^{(2)}(X)\|=\|\mathsf R\|=1$.  Since
$\sqrt2\ne1$, the second amplification of $\alpha$ is not isometric.
\end{proof}

\section{Scope}

The example is finite-dimensional, unital, and nondegenerately represented,
and its acting group is finite.  It therefore gives a full negative answer to
the auxiliary automatic-complete-isometry question at $p=2$.

It does not decide whether the crossed product in the example is
$2$-nuclear exactly when $A$ is $2$-nuclear.  Nor does this row/column proof
extend verbatim to $p\ne2$: the level-one column and row norms then become
the $\ell_p$ and $\ell_q$ norms and are no longer equal.  Hence the main
relaxed theorem and its intended non-Hilbertian range remain outside the
claim of this packet.

\section{Reference}

Z. Wang and S. Zhu, \emph{$p$-nuclearity of $L^p$-operator crossed products},
\href{https://arxiv.org/abs/2305.03933}{arXiv:2305.03933}.

\end{document}

